\section[A uniform hard normal for the q-infinitesimal branch]{A uniform hard normal for the $q$-infinitesimal branch}
This intermediate lemma by itself is not an admission of the full
$q$-infinitesimal C1.1 source. The following two completion sections
supply the scalar allocation and remote joins. Keep (SC1)--(SC3), but do
not discard the whole second paired packet. Its addition to the
$N^2$ coefficient in (SC3) gives exactly
\[
 \frac4R(h_2+h_bD_v)+\frac{4h_bM}{tR}D_v
       =H_0+\frac{4h_b}{t}v^2,
 \qquad H_0=\frac{4h_2}{R}+\frac{4h_b}{t}.
\]
Retain the last summand as the positive square
$4h_b(vN)^2/t$, and use the pure-source normalization
\[
                         \widehat N=\sqrt{H_0}\,N.
\]
The exact original inequality $\beta\ge0$ gives
$h_b/t\ge g^2/2$. The explicit formula for $h_2$ gives
$h_2/q$ a positive bounded unit, uniformly through both openings.
Consequently $H_0\ge c_0(q+g^2)$ for some positive real $c_0$.
This argument also permits $\beta=0$ literally. It does not require
$e=\sqrt q/g$ or its reciprocal to be bounded.

The coefficient of $\widehat N^2$ from this source part is exactly
one. Its remaining contribution $2g^2u^2v^2/H_0$ is bounded and
infinitesimal on the small-angle testing chart. Moreover
\[
 \frac{g}{\sqrt{H_0}},\quad
 \frac{\sqrt q}{\sqrt{H_0}}
\]
are bounded. Since $h_b-h_2=O(qg^2)$, the complete first normal
(SC14), after this normalization, is an infinitesimal-coefficient
linear combination of $B,H,V$ for $u,v$ infinitesimal. The first
coefficient is bounded because
\[
 \frac{h_bD_v-h_2}{\sqrt{qH_0}}
            =O(g^2+v^2),
\]
and the other two coefficients are bounded by constant multiples of
$buv$ and $v$. The hard cross coefficient with $N_y$ is likewise
infinitesimal.

The issue left by this lemma is scalar and angular: $H$ is unweighted in this
statement, whereas the four-term circuit controls $\sqrt q H$.
One must retain a strict $D$-based contact normal and allocate every
unweighted $H$ coefficient on that normal and the circuit. The
naive full extraction (SC10) fails in this branch: its raw $P_H$
satisfies the exact formal specialization
\[
       P_H\big|_{q=0,\,v=0}=-\frac{2ru^2}{1+u^2},
\]
so $P_H/q$ need not be bounded. Thus the corrected hard-normal
normalization removes the former $g/\sqrt q$ obstruction. The later
protected-square lift establishes the missing scalar allocation.
